% 夜半用户手册 · 第 8 章：导出与分享（中文）
% T-DOC-001 逐章交付（CH08 任务未到前预写）。
% 事实来源：app.slint 导出对话框属性、docs/manual/06-export.md、
% export_stems.rs（分轨回执与 manifest）、README.md（CLI 入口）。

\chapter{导出与分享}
\label{ch:export-sharing}

作品写好后要离开夜半：交给别的软件混音、送去母带、或只是留一份可播放的
文件。这一章讲三种出口各导出什么，以及怎么核对导出的结果——「完成」两个字
不等于真的导对了。

\section{WAV 导出}

\subsection{导出整轨混音}

「文件 $\to$ 导出音频 (WAV)...」把整首作品混成一份 WAV 文件。对话框里可以设：

\begin{itemize}
  \item \textbf{输出路径}：可用 Ctrl+O 浏览选择。
  \item \textbf{采样率}：可指定，默认跟随工程（界面显示 \texttt{auto}）。
  \item \textbf{位深}：默认 24 位。
  \item \textbf{导出时长}：以秒计。
  \item \textbf{母线限制器}：默认关闭；打开后渲染会在总线上做峰值限制，防止削波。
\end{itemize}

\screenshot{ch08-export-wav-dialog}{导出音频 (WAV) 对话框，标出路径、采样率、位深、时长与限制器开关}

\subsection{读懂导出回执}

导出回执给出实测读数：帧数、真峰值与限制器的工作量，以及输出位置。这些
数字来自实际渲染，不是从输入框抄回来的；限制器开没开、有没有削波，看回执
而不是凭感觉。

\section{分轨导出（Stems）}

\subsection{什么时候用分轨}

「文件 $\to$ 导出分轨 (Stems)...」把每条源节点轨单独导成一份 WAV，交给别的
软件混音或送母带时用。主总线不出文件，但会在跳过清单里点名——这不是丢轨。

导出在后台线程进行，不卡住界面；完成后状态提示会报出导了几条、跳过几条与
输出目录。每条分轨是 32-bit Float WAV，采样率取自工程；文件名带序号前缀，
并在输出目录里写一份 \texttt{stems-manifest.json} 索引，记下每份文件的轨身份、
文件名、字节数与哈希。

\screenshot{ch08-stems-toast}{分轨导出完成后的状态提示，显示导出条数、跳过条数与输出目录}

\subsection{核对回执，别只看「完成」}

分轨回执以 \texttt{exported-stems:} 开头的一行汇总，随后每条分轨一行、每条
跳过轨一行。除了轨级的「跳过几条」，还应看\textbf{音频片段}维度的渲染/跳过
计数：如果某条分轨里的音频片段全被跳过，提示应降级为警告，而不是仍然只说
「分轨完成」。看到「完成」时，对照回执里的片段计数核对一遍是稳妥习惯。

\begin{figure}[htbp]
\centering
\begin{tikzpicture}[node distance=1.1cm, >=stealth,
    box/.style={draw, rounded corners, align=center, minimum width=2.2cm, minimum height=0.8cm}]
  \node[box] (a) {选出口};
  \node[box, right of=a] (b) {设参数};
  \node[box, right of=b] (c) {后台渲染};
  \node[box, right of=c] (d) {落盘 + manifest};
  \node[box, right of=d] (e) {读回执核对};
  \draw[->] (a)--(b); \draw[->] (b)--(c); \draw[->] (c)--(d); \draw[->] (d)--(e);
\end{tikzpicture}
\caption{导出流程：选出口 $\to$ 设参数 $\to$ 渲染 $\to$ 落盘 $\to$ 读回执核对}
\label{fig:ch08-export-flow}
\end{figure}

命令行同样可以导出分轨：\texttt{yeban-app --open X.yeban --export-stems stems/}；
不开窗口、不建 MCP 会话，回执口径与 GUI 一致。

\section{工程打包与归档}

\texttt{.yeban} 本身就是打包格式：轨道、片段、音符、设备链与所用素材都在一个
容器里，保存是原子写入（见第 2 章）。因此「分享工程」通常就是把这个文件连同
它引用的资产一起交出去；资产池按内容哈希寻址，换机器打开不会变成无声。

「诊断 $\to$ 导出工程诊断报告包」导出的是排障用的诊断包，不是作品音频；作品
有问题要找人看时附上它。

\section{帮助主题映射}

\begin{table}[htbp]
\centering
\caption{本章小节 $\to$ 应用内帮助锚点}
\label{tab:ch08-help}
\begin{tabular}{ll}
\hline
本章小节 & 帮助主题 \\
\hline
导出整轨混音 / 对话框参数 & \texttt{export-wav-dialog} \\
读懂导出回执 & \texttt{export-wav-receipt} \\
什么时候用分轨 & \texttt{export-stems-purpose} \\
核对回执，别只看「完成」 & \texttt{export-stems-receipt} \\
工程打包与归档 & \texttt{export-project-archive} \\
\hline
\end{tabular}
\end{table}

\paragraph{本章学完能做什么} 按用途选对导出方式，设置对话框参数，导出后用回执（含片段计数）核对产物，并知道诊断包不是作品。

\paragraph{动手练习} 先导出整轨 WAV，打开限制器，读回执里的帧数与真峰值；再导出分轨，对照 \texttt{stems-manifest.json} 核对文件名与字节数，并从回执里找出音频片段的渲染/跳过计数；最后用命令行对同一工程导出一次分轨，比较两路音频样本是否一致。
